\chapter{Conclusions and Future Work}\label{ch:conclusions}

\section{Conclusions}\label{sec:conc-conclusion}
This thesis started from a simple observation about tracking-by-detection systems: the settings that connect the detector and the tracker---the detector operating point and the score thresholds of the tracker---are usually fixed once for one detector and one kind of scene, although the scene changes within a video and the detector changes during the life of a system. The research problem was to adapt these settings online, from information that is available while the video is processed, without labels, without training, and with the detector and the tracker left unchanged (Section~\ref{sec:intro-problem}).

The published survey \cite{ismail2026realtime} provided the starting point. It reviewed the detectors, trackers, datasets, and metrics of real-time detection and tracking, showed that tracking by detection remains the dominant paradigm, and pointed to the dependence between detection accuracy and association reliability under real-time constraints. Together with the review of related work in Chapter~\ref{ch:background}, it identified the gap addressed by the thesis: existing adaptive methods either learn a policy for one detector, adapt one tracker from inside, or require labeled data.

The first controller, AC-MOT (Chapter~\ref{ch:acmot}), answered RQ1 for UAV video. It adapts the confidence threshold, the NMS threshold, and the input resolution of the detector every ten frames. A scene complexity index formed from the crowd cue, checks from four further inexpensive scene cues, and a three-state controller with recovery, hysteresis, and dwell rules choose one of three operating points; the values of the controller were selected by cross-validation on separate calibration sequences and then frozen. On seven VisDrone2019-MOT sequences it improved HOTA by 2.937 points with ByteTrack and 1.325 points with OATrack, and MOTA by 17.995 and 6.318 points, with fewer identity switches and false positives; paired bootstrap analysis supports these gains. On UAVDT it improved MOTA by 17.601 and 3.267 points without retuning, and it adds about 4.2 ms per frame on a Tesla T4.

The second controller, SCL (Chapter~\ref{ch:scl}), answered RQ2. It reads the score distribution of recent frames with nested Otsu splits and a whole-stream majority rule. On clean scores, it keeps the tracker's scores and, where they fit, its thresholds; on noisy scores, it chooses the candidates, rescores them, and sets the thresholds. A track-guided recovery gives single-stage trackers the continuation step they lack, and a host contract of four published settings lets one configuration serve different trackers. One frozen configuration restored two trackers from 0 to about 66 HOTA under halved scores, improved both single-stage trackers, including one not seen during design, on every MOT17 sequence, raised MOTA by up to 24.70 points with a noisy detector on held-out aerial data, cut the identity switches of two-stage trackers by 43\% to 56\% with replacement detectors, and kept well-matched trackers at their own accuracy. It reached higher HOTA than simple adaptive thresholds (by 2.2 to 8.6 points) and than the published ACT method, and it takes 2.95 ms per frame on an ordinary processor.

Both controllers kept their effect on data, trackers, or detectors that were not used to design them, at a cost of a few milliseconds per frame, which answers RQ3 for the tested settings: AC-MOT on UAVDT without retuning, and SCL on two predeclared external trackers and on the held-out UAVDT test split. Table~\ref{tab:conc-summary} summarizes the two controllers and their main results, and Table~\ref{tab:conc-rq} relates each research question to its principal evidence and main finding.

\begin{table}[tbp]
\centering
\caption{Summary of the two proposed controllers and their main results; all changes are relative to the same pipeline without the controller.}
\label{tab:conc-summary}
\begin{tabular}{>{\raggedright\arraybackslash}p{2.7cm}>{\raggedright\arraybackslash}p{5.6cm}>{\raggedright\arraybackslash}p{5.6cm}}
\toprule
Aspect & AC-MOT & SCL\\
\midrule
Question & RQ1: adapt the detector operating point to the scene & RQ2: restore the compatibility of score scale and tracker thresholds\\
Acts on & detector operating point & tracker input and thresholds\\
Evaluated on & VisDrone2019-MOT and UAVDT with ByteTrack and OATrack & MOT17, KITTI, and UAVDT with seven published trackers\\
Main results & VisDrone: HOTA $+2.937$ and $+1.325$, MOTA $+17.995$ and $+6.318$; UAVDT: MOTA $+17.601$ and $+3.267$ & halved scores: 0 to 66.17 and 66.72 HOTA; KITTI: MOTA $+5.98$ and $+7.86$; UAVDT: MOTA $+24.70$\\
Cost & 4.26 and 4.20 ms per frame (Tesla T4) & 2.95 ms per frame (CPU)\\
\bottomrule
\end{tabular}
\end{table}

\begin{table}[tbp]
\centering
\caption{Research questions of the thesis, the principal experiments that answer them, and the main findings.}
\label{tab:conc-rq}
\begin{tabular}{>{\raggedright\arraybackslash}p{2.9cm}>{\raggedright\arraybackslash}p{5.3cm}>{\raggedright\arraybackslash}p{5.7cm}}
\toprule
Research question & Principal evidence & Main finding\\
\midrule
RQ1: online control of the detector operating point & AC-MOT on seven VisDrone2019-MOT sequences with ByteTrack and OATrack, and frozen transfer to UAVDT (Sections~\ref{sec:ac-res-visdrone}--\ref{sec:ac-res-uavdt}) & over the fixed operating point: HOTA $+2.937$ and $+1.325$ and MOTA $+17.995$ and $+6.318$ on VisDrone; MOTA $+17.601$ and $+3.267$ on UAVDT without retuning\\
RQ2: label-free compatibility of scores and tracker thresholds & SCL under halved scores, with replacement detectors, and on MOT17, KITTI, and UAVDT, compared with adaptive thresholds and ACT (Sections~\ref{sec:scl-results} and~\ref{sec:scl-compare}) & two trackers restored from 0 to about 66 HOTA; well-matched trackers kept at their own accuracy; higher HOTA than simple adaptive thresholds (by 2.2 to 8.6 points) and than ACT\\
RQ3: effect on unseen data, trackers, and detectors, and cost & AC-MOT on UAVDT without retuning; SCL on two predeclared external trackers and the held-out UAVDT test split; runtime (Sections~\ref{sec:ac-res-uavdt}, \ref{sec:ac-res-runtime}, \ref{sec:scl-res-external}, \ref{sec:scl-res-uavdt}, and~\ref{sec:scl-runtime}) & effects kept on data, trackers, and detectors not used in design; 4.2 ms per frame for AC-MOT (Tesla T4) and 2.95 ms per frame for SCL (CPU)\\
\bottomrule
\end{tabular}
\end{table}

Taken together, the two studies advance the understanding of robust tracking in two ways. First, they show that the operating point of the detector and the score scale that the tracker receives are parts of the tracking system that deserve the same attention as the association method: with the detector and the tracker unchanged, each of them moved MOTA or HOTA by many points. Second, they show that both can be controlled from outside the detector and the tracker, with rules whose values were fixed before the final evaluation: AC-MOT selects the detector settings that suit the current aerial scene, and SCL makes the scores of a given detector usable by different trackers and leaves a well-matched pair unchanged.

\emph{Practical implications.} Neither controller needs training, labels, or access to the inside of the detector or the tracker. AC-MOT changes three arguments of the detector call, and SCL sets the published thresholds and the input of the tracker for each frame. A system integrator can therefore keep a tuned tracker when the scenes change or when the detector is replaced or retrained, and the label-free mode of SCL can be logged as a running check that a new detector still fits its tracker.

\emph{Evaluation coverage.} AC-MOT was evaluated on two aerial datasets, VisDrone2019-MOT and UAVDT, with two trackers of different design, ByteTrack and OATrack, and improved MOTA on UAVDT without retuning. SCL was evaluated with one frozen configuration on seven published tracking-by-detection trackers, four detection sources, and three datasets, including two predeclared external trackers and a held-out test split. The two controllers were developed and evaluated separately; their joint use in one pipeline was not part of the evaluation.

\section{Future Work}\label{sec:conc-future}
The training-free and modular design of the two controllers opens several directions that build on these results:
\begin{enumerate}
\item \emph{New detectors and aerial platforms.} Because AC-MOT needs no training and changes only three arguments of the detector call, the same calibration procedure can be applied to new detector families and aerial datasets, and its action map can include further operating points. Each new detector is to be evaluated in two separate conditions, the frozen controller transferred without retuning and the controller recalibrated with the same cross-validation procedure, by HOTA, MOTA, and IDF1 against the detector at its fixed operating point.
\item \emph{End-to-end trackers.} The host contract of SCL can be extended to trackers that learn detection and association jointly, so that the same layer also serves this family of trackers; as in Chapter~\ref{ch:scl}, the objective is judged by HOTA and IDS against the same tracker alone under score shifts and detector changes.
\item \emph{Monitoring of detector updates.} The label-free mode and the confident share of SCL can serve as a running check when the detector of a deployed system is updated, retrained, or compressed; the objective is to log the mode and $\rho$ before and after each update and to relate them to the tracking accuracy measured on labeled test sequences.
\item \emph{A combined pipeline.} Since the two controllers act on different parts of the pipeline, they can be run together, with AC-MOT setting the detector operating point and SCL adapting the tracker input to the resulting scores, to evaluate whether their improvements are complementary. The planned evaluation compares four configurations on the same sequences---the original pipeline, AC-MOT alone, SCL alone, and both controllers together---by tracking accuracy (HOTA and MOTA), identity preservation (IDF1 and IDS), the stability of the controller states and of the SCL modes (changes per 100 frames), and the processing time per frame.
\item \emph{Onboard deployment.} Both controllers add only a few milliseconds per frame, which makes an implementation on the processors used on board UAVs a natural next step. The target is an embedded GPU of the kind carried on UAVs, such as an NVIDIA Jetson module, with the whole pipeline at 30 FPS, that is, at most 33.3 ms per frame for the detector, the controllers, and the tracker together, measured with the power use and the memory.
\end{enumerate}
